\section{Seudocódigo}

El seudocódigo es básicamente algo que se parece a un lenguaje de programación pero que permite:
\begin{itemize}
  \item Despreocuparse de detalles de sintaxis o de implementación, como el manejo de errores.
  \item Usar la forma más clara de expresarse en cada línea: a veces español, a veces matemática, a veces sintaxis de algún lenguaje de programación.
\end{itemize}
Entonces es muy útil para expresar algoritmos.

En los textos formales se definen convenciones estrictas para el seudocódigo que se usa; eso elimina la ambigüedad, pero obliga al lector a aprenderse las convenciones, prácticamente un nuevo lenguaje. Contrariamente, en estos apuntes se utiliza un seudocódigo menos estructurado, basado en la sintaxis de Python, pero con varias libertades. Las ambigüedades se eliminan aclarándolas cuando aparecen.

En Python son muy idiomáticas la instrucciones de alto nivel, que automáticamente ordenan, buscan, etc. y que no existen en otros lenguajes. Ese tipo de instrucciones se omiten deliberadamente, primero porque formalmente arruinarían el análisis algorítmico, como se expone más adelante, y segundo porque la idea es precisamente entender cómo implementar eficientemente ese tipo de instrucciones. Así que, bajo esa perspectiva, el seudocódigo es más parecido a un lenguaje de programación de mediano o bajo nivel.

E.g., una función que recibe un arreglo de números y encuentra el mayor:
\begin{pseudocode}
función encontrar_mayor(números: arreglo de enteros) entero
    mayor = |$-\infty$|
    for i=0 to números.longitud-1
        if números[i] > mayor
            mayor = números[i]
    return mayor
\end{pseudocode}
Detalles de este seudocódigo que reflejan que su propósito es ser claro:
\begin{itemize}
  \item Algunas palabras y tipos están en español. A pesar de que no se definieron de antemano, son muy claros porque son los términos que usamos para referirnos a esas cosas.
  \item Omito detalles de sintaxis como los dos puntos, pero preservo los que otorgan claridad, como la indentación.
  \item Evito sintaxis específica de algún lenguaje.
  \begin{itemize}
    \item En el for uso una sintaxis que hace explícito el valor inicial y el valor final sin necesidad de usar iterables (y sin tener que desambibuar si se toma o no el útlimo valor; acá es claro que sí)
    \item Presumo la existencia de un método \texttt{longitud} para obtener la longitud de un arreglo. Es una presunción razonable porque en cualquier lenguaje orientado a objetos existe algo así.
  \end{itemize}
\end{itemize}